% Standalone диаграмма архитектуры
% xelatex arch_diagram.tex
\documentclass[border=12pt, right=15pt]{standalone}

\usepackage{fontspec}
\usepackage{polyglossia}
\setdefaultlanguage{russian}
\setmainfont{DejaVu Sans}
\setsansfont{DejaVu Sans}

\usepackage{tikz}
\usetikzlibrary{arrows.meta, positioning, fit, backgrounds, shadows.blur, calc}

\definecolor{CTg}   {RGB}{0,   136, 204}
\definecolor{CMst}  {RGB}{25,  50,  105}
\definecolor{CMem}  {RGB}{10,  105,  80}
\definecolor{CLrn}  {RGB}{140,  60,   0}
\definecolor{CDiag} {RGB}{75,  25, 130}
\definecolor{CAut}  {RGB}{155,  15,  15}
\definecolor{CGen}  {RGB}{35,   90, 160}
\definecolor{CNew}  {RGB}{20,  120,  50}
\definecolor{BGcore}{RGB}{238, 244, 255}
\definecolor{BGauto}{RGB}{255, 240, 240}
\definecolor{Clbl}  {RGB}{80,   85, 105}
\definecolor{Cwht}  {RGB}{252, 253, 255}

\tikzset{
  N/.style={
    draw, rounded corners=4pt,
    text=white, align=center,
    font=\fontsize{8}{9.5}\selectfont,
    minimum width=32mm, minimum height=10mm,
    blur shadow={shadow blur steps=3,
                 shadow xshift=0.5mm, shadow yshift=-0.5mm,
                 shadow opacity=18}
  },
  Ntg/.style={N, fill=CTg,   draw=CTg!80!black,
               minimum width=36mm,
               font=\fontsize{8.5}{10}\bfseries},
  Nmst/.style={N, fill=CMst,  draw=CMst!70!black,
               minimum width=36mm, minimum height=11mm,
               font=\fontsize{8.5}{10}\bfseries},
  Nmem/.style={N, fill=CMem,  draw=CMem!70!black},
  Nlrn/.style={N, fill=CLrn,  draw=CLrn!70!black},
  Ndiag/.style={N, fill=CDiag, draw=CDiag!70!black, minimum width=36mm},
  Naut/.style={N, fill=CAut,  draw=CAut!70!black,  minimum width=32mm},
  Ngen/.style={N, fill=CGen,  draw=CGen!70!black},
  Nnew/.style={N, fill=CNew,  draw=CNew!70!black,  minimum width=32mm},
  A/.style={-{Stealth[length=4pt,width=3.5pt]}, line width=0.85pt},
  Ad/.style={A, dashed},
  L/.style={font=\fontsize{6.5}{7.5}\selectfont,
             fill=Cwht, inner sep=1.5pt, text=Clbl, rounded corners=1.5pt},
}

\begin{document}
\begin{tikzpicture}[x=1mm, y=1mm]

% ── Узлы ─────────────────────────────────────────────────────────
\node[Ntg]   at (  0,   0)  (tg)  {Telegram};
\node[Nmst]  at (  0, -22)  (mst) {Мастер (SBCL)};

\node[Nlrn]  at (-52, -44)  (lrn) {поток-обучение};
\node[Nmem]  at (  0, -44)  (mem) {поток-память};
\node[Ndiag] at ( 56, -44)  (dia) {поток-диагностика};

\node[Naut]  at (  0, -66)  (aut) {поток-автоматор};
\node[Ngen]  at ( 56, -66)  (gen) {порождатель};
\node[Nnew]  at ( 56, -86)  (nw)  {новый поток};

% ── Подложки ─────────────────────────────────────────────────────
\begin{scope}[on background layer]
  \node[fill=BGcore, draw=CMst!15, rounded corners=7pt,
        inner sep=7pt, fit=(lrn)(mem)(dia),
        label={[font=\fontsize{6.5}{8}\selectfont,
                text=CMst!45!black,
                anchor=north west, xshift=4pt, yshift=2pt]north west:потоки ядра}
  ] (grpcore) {};
  \node[fill=BGauto, draw=CAut!15, rounded corners=7pt,
        inner sep=7pt, fit=(aut)(gen)(nw),
        label={[font=\fontsize{6.5}{8}\selectfont,
                text=CAut!55!black,
                anchor=north west, xshift=4pt, yshift=2pt]north west:автоматизация}
  ] (grpauto) {};
\end{scope}

% ── Вспомогательные точки ────────────────────────────────────────
\coordinate (mst_bot) at ($(mst.south)+(0,-6)$);

% ── Стрелки ──────────────────────────────────────────────────────

% Telegram → Мастер
\draw[A, color=CTg!75!black]
  (tg.south) -- (mst.north)
  node[L, midway, right=4pt] {сообщения};

% Мастер → обучение  (вниз → влево → вниз)
\coordinate (lrn_top) at ($(lrn.north)+(0,0)$);
\draw[A, color=CLrn!75!black]
  (mst.south west) -- ++(0,-6) -- (-52,-28) -- (lrn.north)
  node[L, pos=0.6, above=1.5pt] {/обучить};

% Мастер → память  (прямо вниз)
\draw[A, color=CMem!75!black]
  (mst.south) -- (mem.north)
  node[L, midway, right=3pt] {/список};

% Мастер → диагностика  (вниз → вправо → вниз)
\draw[A, color=CDiag!75!black]
  (mst.south east) -- ++(0,-6) -- (56,-28) -- (dia.north)
  node[L, pos=0.6, above=1.5pt] {/диагностика};

% обучение → память  (горизонтально, над центром)
\draw[A, color=CLrn!80!black]
  (lrn.east) -- (mem.west)
  node[L, midway, above=2pt] {сохранить};

% диагностика → память  (горизонтально, ниже центра, пунктир)
\draw[Ad, color=CMem!75!black]
  (dia.west) -- (mem.east)
  node[L, midway, below=2pt] {найти похожие};

% память → автоматор  (вниз)
\draw[A, color=CMem!70!black]
  (mem.south) -- (aut.north)
  node[L, midway, right=3pt] {паттерны};

% диагностика → автоматор  (вниз → влево → авт.восток)
\coordinate (d_below) at ($(dia.south)+(0,-12)$);
\coordinate (a_east)  at ($(aut.east)+(0,0)$);
\draw[Ad, color=CDiag!55!black]
  (dia.south) -- (d_below) -- ($(a_east)+(0,-12)$) -- (aut.east)
  node[L, pos=0.45, right=2pt] {уверенность};

% автоматор → порождатель  (горизонтально)
\draw[A, color=CAut!80!black]
  (aut.east) -- (gen.west)
  node[L, midway, above=2pt] {описание};

% порождатель → новый поток  (вниз)
\draw[A, color=CGen!80!black]
  (gen.south) -- (nw.north)
  node[L, midway, right=3pt] {compile+load};

% новый поток → Мастер  (правая шина: вправо → вверх → влево)
\coordinate (bus_x) at (76,-86);
\draw[Ad, color=CNew!65!black, rounded corners=3pt]
  (nw.east) -- ++(8,0) -- ++(0,75) -- (mst.east)
  node[L, pos=0.45, right=3pt] {→ в образ};

% feedback: Мастер → обучение  (левая шина)
\draw[Ad, color=CLrn!50!black, rounded corners=3pt]
  (mst.west) -- ++(-8,0) -- ++(0,-22) -- (lrn.north west)
  node[L, pos=0.45, left=3pt] {feedback};

\end{tikzpicture}
\end{document}
